\subsection{Symbols}

14.2.1. Let \(k\) be a positive integer \(\leqslant r\) . Consider the exact sequence

\[
0 \rightarrow \mathrm{P} _ {k} (\mathrm{T} (\mathrm{X}); \mathrm{E}) \xrightarrow {\iota} \mathrm{P} ^ {k} (\mathrm{E}) \xrightarrow {\jmath} \mathrm{P} ^ {k - 1} (\mathrm{E}) \rightarrow 0
\]

of no. 12.6.8, where we have set \(j = r^{k, k-1}\) . Applying the vector functor M → L(M; F) to this sequence, we obtain the exact sequence

\[
0 \rightarrow \mathscr {L} (\mathrm{P} ^ {k - 1} (\mathrm{E}); \mathrm{F}) \xrightarrow {f ^ {\prime}} \mathscr {L} (\mathrm{P} ^ {k} (\mathrm{E}); \mathrm{F}) \xrightarrow {\sigma_ {k}} \mathscr {L} (\mathrm{P} _ {k} (\mathrm{T} (\mathrm{X}); \mathrm{E}); \mathrm{F}) \rightarrow 0
\]

which can be written:

by setting

\[
0 \rightarrow \mathrm{D} ^ {k - 1} (\mathrm{E}, \mathrm{F}) \xrightarrow {j ^ {\prime}} \mathrm{D} ^ {k} (\mathrm{E}, \mathrm{F}) \xrightarrow {\sigma_ {k}} \mathrm{S} ^ {k} (\mathrm{E}, \mathrm{F}) \rightarrow 0
\]

\[
\mathbf {S} ^ {k} (\mathrm{E}, \mathrm{F}) = \mathcal {L} (\mathrm{P} _ {k} (\mathrm{T} (\mathrm{X}); \mathrm{E}); \mathrm{F}).
\]

The homomorphism \(j' = \mathcal{L}(j; \mathrm{Id}_{\mathrm{F}})\) is the canonical inclusion of \(\mathbf{D}^{k-1}(\mathbf{E}, \mathbf{F})\) into \(\mathbf{D}^k(\mathbf{E}, \mathbf{F})\) ; the homomorphism \(\sigma_k\) is by definition equal to \(\mathcal{L}(\iota; \mathrm{Id}_{\mathrm{F}})\) . If D is a differential operator of type \(\mathbf{E} \to \mathbf{F}\) and of order \(\leqslant k\) , its image under \(\sigma_k\) is a section \(\sigma_k(\mathbf{D})\) of \(\mathbf{S}^k(\mathbf{E}, \mathbf{F})\) which we call the \(k\)-symbol (or simply the symbol) of D; we have \(\sigma_k(\mathbf{D}) = 0\) if and only if D is of order \(\leqslant k - 1\) . If D is of class \(\mathbf{C}^h\) so is its symbol.

14.2.2 (“Identifications of the symbol bundle”). One may write the bundle \(S^{k}(E, F) = \mathcal{L}(P_{k}(T(X); E); F)\) in various ways. First, if \(x \in X\) , an element of \(P_{k}(T_{x}(X); E_{x})\) is identified (13.1.2) with an element of \(\mathcal{L}(\mathsf{T}\mathsf{S}^{k}(\mathsf{T}_{x}(X)); E_{x})\) . We thus obtain identifications of vector bundles

\[
\mathrm{P} _ {k} (\mathrm{T} (\mathrm{X}); \mathrm{E}) = \mathscr {L} (\mathrm{TS} ^ {k} (\mathrm{T} (\mathrm{X})); \mathrm{E}) = (\mathrm{TS} ^ {k} (\mathrm{T} (\mathrm{X}))) ^ {*} \otimes \mathrm{E}
\]

and, by applying the vector functor \(\mathbf{M} \mapsto \mathcal{L}(\mathbf{M}; \mathbf{F}) = \mathbf{M}^* \otimes \mathbf{F}\) we obtain:

\[
\mathrm{S} ^ {k} (\mathrm{E}, \mathrm{F}) = \mathrm{T} \mathrm{S} ^ {k} (\mathrm{T} (\mathrm{X})) \otimes \mathrm{E} ^ {*} \otimes \mathrm{F} = \mathrm{T} \mathrm{S} ^ {k} (\mathrm{T} (\mathrm{X})) \otimes \mathscr {L} (\mathrm{E}; \mathrm{F}).
\]

One can further transform the preceding expression. If M and N are two vector bundles, denote by \(\operatorname{Sym}^{k}(\mathbf{M};\mathbf{N})\) the subbundle of \(\mathcal{L}(\mathbf{M},\ldots,\mathbf{M};\mathbf{N})\) consisting of symmetric \(k\)-linear maps from \(\mathbf{M}^{k}\) to \(\mathbf{N}\). There is a canonical isomorphism

\[
\operatorname{Sym} ^ {k} (\mathbf {M}; \mathbf {N}) \rightarrow \operatorname{Sym} ^ {k} (\mathbf {M}; \mathbf {K} _ {\mathbf {X}}) \otimes \mathbf {N};
\]

On the other hand, the duality between \(\otimes^{k}\mathbf{M}\) and \(\otimes^{k}\mathbf{M}^{*}\) identifies \(\operatorname{Sym}^{k}(\mathbf{M};\mathbf{K}_{\mathbf{x}})\) with \(\operatorname{TS}^{k}(\mathbf{M}^{*})\) . One thus obtains an identification

\[
\operatorname{Sym} ^ {k} (\mathbf {M}; \mathbf {N}) = \operatorname{TS} ^ {k} (\mathbf {M} ^ {*}) \otimes \mathbf {N}.
\]

Applying this formula to M = T(X)* and N = ℒ(E; F), we have:

\[
\mathrm{S} ^ {k} (\mathrm{E}, \mathrm{F}) = \mathrm{T} \mathrm{S} ^ {k} (\mathrm{T} (\mathrm{X})) \otimes \mathscr {L} (\mathrm{E}; \mathrm{F}) = \operatorname{Sym} ^ {k} (\mathrm{T} (\mathrm{X}) ^ {*}; \mathscr {L} (\mathrm{E}; \mathrm{F})).
\]

Suppose K has characteristic 0 or \(>k\). If u is a symmetric \(k\)-linear map, denote \(\tilde{u}\) the polynomial map

\[
x \mapsto \frac {1}{k !} u (x, \dots , x)
\]

This is the corresponding polynomial map. The map \(u \mapsto \tilde{u}\) defines an isomorphism

\[
\mathrm{S} ^ {k} (\mathrm{E}, \mathrm{F}) = \operatorname{Sym} ^ {k} (\mathrm{T} (\mathrm{X}) ^ {*}, \mathscr {L} (\mathrm{E}; \mathrm{F})) \rightarrow \mathrm{P} _ {k} (\mathrm{T} (\mathrm{X}) ^ {*}; \mathscr {L} (\mathrm{E}; \mathrm{F})).
\]

Thus, every element \(\sigma\) of \(S^{k}(E,F)_{x}\) , with \(x\in X\), can be identified with a homogeneous polynomial map \(\tilde{\sigma}\) of degree \(k\) from \(T_{x}(X)^{*}\) to \(\mathcal{L}(E_{x};F_{x})\) .

14.2.3 (“Computation of the symbol of a differential operator”). Let \(x \in X\) and let D be a differential operator of type \(E \to F\) of order \(\leqslant k\) defined in an open neighborhood of \(x\) and of class \(\mathbf{C}^{r-k}\) . We will now make explicit the value \(\sigma_k(D)(x)\) of the symbol of D at \(x\) . Consider it as an element of \(\text{Sym}^k(T_x(X)^*; \mathcal{L}(E_x; F_x))\) (cf. 14.2.2). Let \(\omega_1, \ldots, \omega_k\) be covectors at \(x\), and choose functions \(f_1, \ldots, f_k\) of class \(\mathbf{C}^r\) in an open neighborhood U of \(x\) such that \(d_x f_i = \omega_i\) for \(i = 1, \ldots, k\). Set

\[
\mathbf {D} ^ {\prime} = (- 1) ^ {k} \mathrm{ad} (f _ {k}) \dots \mathrm{ad} (f _ {1}) \mathbf {D} \quad (\text { cf. } 14.1.10, a).
\]

The operator D is of order \(\leqslant0\); it is a section of \(\mathcal{L}(\mathrm{E};\mathrm{F})|\mathrm{U}\) . The value of \(\mathbf{D}'\) at x is an element \(\lambda(\omega_{1},\ldots,\omega_{k})\) of \(\mathcal{L}(\mathrm{E}_{x};\mathrm{F}_{x})\) which depends only on \((\omega_{1},\ldots,\omega_{k})\) . The map \(\lambda\) thus defined is symmetric. It is equal to the symbol \(\sigma_{k}(\mathbf{D})(x)\) of D at x.

Suppose K = R or C, and let us make explicit \(\sigma_{k}(\mathbf{D})(x)\) considered as a homogeneous polynomial map of degree k from \(\mathrm{T}_{x}(\mathrm{X})^{*}\) to \(\mathcal{L}(\mathrm{E}_{x};\mathrm{F}_{x})\) (cf. 14.2.2). Let \(\omega \in \mathrm{T}_{x}(\mathrm{X})^{*}\) and \(v \in E_{x}\) ; choose a function f of class \(C^{r}\) on an open neighborhood U of x such that \(d_{x}f = \omega\) , and a section s of class \(C^{r}\) of E over U such that \(s(x) = v\) . There exists a family \((\varphi_{0}, \varphi_{1}, \ldots, \varphi_{k})\) of sections of F over U, and only one, such that

\[
e ^ {- t f} \mathrm{D} (e ^ {t f} s) = \sum_ {j = 0} ^ {k} t ^ {j} \varphi_ {j} \quad \text { for all } \quad t \in \mathbf {K}.
\]

Then \(\varphi_{k}(x)\) does not depend on the choice of \(f\) and \(s\) such that \(d_{x}f = \omega\) , \(v(x) = s\) ; the map \(v \mapsto \varphi_{k}(x)\) is an element \(\lambda(\omega)\) of \(\mathcal{L}(\mathrm{E}_{x}; \mathrm{F}_{x})\) , and \(\lambda\) is a homogeneous polynomial map of degree \(k\) from \(\mathrm{T}_{x}(\mathrm{X})^{*}\) to \(\mathcal{L}(\mathrm{E}_{x}; \mathrm{F}_{x})\) . This map is equal to the symbol \(\sigma_{k}(\mathrm{D})(x)\) of D at \(x\) .

14.2.4 ("The scalar case"). Take \(\mathbf{E} = \mathbf{F} = \mathbf{K}_{\mathbf{x}}\) , in which case one identifies \(\mathbf{D}^k(\mathbf{E}, \mathbf{F})\) with the bundle \(\mathbf{T}^{(k)}(\mathbf{X})\) of point distributions of order \(\leqslant k\) , cf. 14.1.6. One has

\[
\mathrm{S} ^ {k} (\mathrm{E}, \mathrm{F}) = \mathrm{TS} ^ {k} (\mathrm{T} (\mathrm{X}))
\]

and the symbol

\[
\sigma_ {k} \colon \mathrm{T} ^ {(k)} (\mathrm{X}) \rightarrow \mathrm{TS} ^ {k} (\mathrm{T} (\mathrm{X}))
\]

is none other than the composite

\[
\mathbf {T} ^ {(k)} (\mathbf {X}) \rightarrow \mathbf {T} ^ {(k)} (\mathbf {X}) / \mathbf {T} ^ {(k - 1)} (\mathbf {X}) \xrightarrow {i _ {k}} \mathsf {T S} ^ {k} (\mathbf {T} (\mathbf {X})),
\]

where \(i_{k}\) is the isomorphism defined in 13.3.2.

Take for example X an open subset of \(K^{n}\) , and let \(\{e_{1},\ldots,e_{n}\}\) be the canonical basis of \(K^{n}\) (identified with the tangent spaces \(T_{x}(X)\) ). If \(\alpha\in N^{n}\) is such that \(|\alpha|\leqslant k\) , the symbol of the operator \(\Delta^{\alpha}\) is given by the formulas (cf. 13.3.3):

\[
\begin{array}{l l} \sigma_ {k} (\Delta^ {\alpha}) (x) = 0 & \text { if } \quad | \alpha | <   k \\ \sigma_ {k} (\Delta^ {\alpha}) (x) = \gamma_ {\alpha_ {1}} (e _ {1}) \ldots \gamma_ {\alpha_ {n}} (e _ {n}) & \text { if } \quad | \alpha | = k, \end{array}
\]

the product of the \(\gamma_{\alpha_i}(e_i)\) being the symmetric product, cf. 13.2.6 and 13.3.3.

14.2.5 ("Symbol of a composite"). The notations and hypotheses are those of 14.1.8. The composition of linear maps defines a pairing

\[
\mathcal {L} (\mathrm{F}; \mathrm{G}) \times \mathcal {L} (\mathrm{E}; \mathrm{F}) \rightarrow \mathcal {L} (\mathrm{E}; \mathrm{G}).
\]

On the other hand, the operation of symmetric product defines a pairing

\[
\mathsf {T S} ^ {k ^ {\prime \prime}} (\mathbf {T} (\mathbf {X})) \times \mathsf {T S} ^ {k ^ {\prime}} (\mathbf {T} (\mathbf {X})) \rightarrow \mathsf {T S} ^ {k} (\mathbf {T} (\mathbf {X})).
\]

Since one has

\[
\left. \begin{array}{l} \mathrm{S} ^ {k ^ {\prime \prime}} (\mathrm{F}, \mathrm{G}) = \mathrm{TS} ^ {k ^ {\prime \prime}} (\mathrm{T} (\mathrm{X})) \otimes \mathscr {L} (\mathrm{F}; \mathrm{G}) \\ \mathrm{S} ^ {k ^ {\prime}} (\mathrm{E}, \mathrm{F}) = \mathrm{TS} ^ {k ^ {\prime}} (\mathrm{T} (\mathrm{X})) \otimes \mathscr {L} (\mathrm{E}; \mathrm{F}) \\ \mathrm{S} ^ {k} (\mathrm{E}, \mathrm{G}) = \mathrm{TS} ^ {k} (\mathrm{T} (\mathrm{X})) \otimes \mathscr {L} (\mathrm{E}; \mathrm{G}) \end{array} \right\} \quad \text { cf. } 1 4. 2. 2,
\]

one deduces from this, by tensor product, a pairing

\[
\mathrm{S} ^ {k ^ {\prime \prime}} (\mathrm{F}, \mathrm{G}) \times \mathrm{S} ^ {k ^ {\prime}} (\mathrm{E}, \mathrm{F}) \rightarrow \mathrm{S} ^ {k} (\mathrm{E}, \mathrm{G}).\tag{*}
\]

One then has the formula

\[
\sigma_ {k} (\mathbf {D} ^ {\prime \prime} \circ \mathbf {D} ^ {\prime}) = \sigma_ {k ^ {\prime \prime}} (\mathbf {D} ^ {\prime \prime}). \sigma_ {k ^ {\prime}} (\mathbf {D} ^ {\prime}),
\]

where the product appearing in the right-hand term is defined by the pairing (*) above.

Now suppose K is of characteristic zero, and let \(x \in X\) . Denote by \(\sigma\) (resp. \(\sigma'\) , \(\sigma''\) ) the k-symbol (resp. the \(k'\) -symbol, the \(k''\) -symbol) of \(D'' \circ D'\) (resp. \(D'\) , \(D''\) ) at x. For every \(\omega \in \mathbf{T}_{x}(\mathbf{X})^{*}\) , one then has (with the notations of the end of 14.2.2):

\[
\tilde {\sigma} (\omega) \in \mathscr {L} (\mathrm{E} _ {x}; \mathrm{G} _ {x}), \quad \tilde {\sigma} ^ {\prime} (\omega) \in \mathscr {L} (\mathrm{E} _ {x}; \mathrm{F} _ {x}), \quad \tilde {\sigma} ^ {\prime \prime} (\omega) \in \mathscr {L} (\mathrm{F} _ {x}; \mathrm{G} _ {x})
\]

and

\[
\tilde {\sigma} (\omega) = \tilde {\sigma} ^ {\prime \prime} (\omega) \circ \tilde {\sigma} ^ {\prime} (\omega).
\]

